% Standalone section material; the parent article supplies theorem environments.
\section{First-jet decisions before cobordism elimination}
\label{sec:first-jet-pre-schur}

The maintained closure observer reads a radical complex after all scalar
units have been cancelled. It deliberately declines a block containing a
scalar coefficient. This section moves the observation before that operation
and proves a rigidity statement that removes the knot-completion hypothesis
when the multiplier is one. The result is an optional decision scanner;
no original graded homology is claimed after a replacement.

Fix a nonempty matching $m$ with $p$ arcs and put
\[
 R=\mathbb F_2[x_1,\ldots,x_p]/(x_1^2,\ldots,x_p^2),\qquad
 A_\epsilon=\mathbb F_2[\epsilon]/(\epsilon^2).
\]
Let $C=(R\otimes V,Q)$ be a bounded finite free complex supported entirely
on $m$. It must be a whole relative direct summand: no omitted differential
entry may attach it to another object. Write
\[
 Q=A+\sum_{i=1}^p x_i B_i+\text{terms of dot degree at least two},
 \qquad B=\sum_i B_i,\qquad N=\dim V .
\]
All matrices raise homological degree by one. The identities $A^2=0$ and
$AB+BA=0$ follow from $Q^2=0$. They are necessary first-jet tests, not
sufficient tests of the whole polynomial identity $Q^2=0$.

\subsection{Two ordinary binary ranks}

Let $r(M)$ denote the total rank of a matrix on the direct sum of all
homological degrees, equivalently the sum of its adjacent-degree ranks.
Define
\begin{equation}
 t=N-2r(A),\qquad
 \kappa=N-r\begin{pmatrix}A&0\\ B&A\end{pmatrix}.
 \label{eq:first-jet-two-ranks}
\end{equation}

\begin{proposition}[Pre-elimination first-jet invariants]
\label{prop:first-jet-invariants}
The number $t$ is the number of objects in a completely scalar-reduced
one-matching complex. If $B_*$ is induced by $B$ on $H(V,A)$, then
\[
 \kappa=t-r(B_*).
\]
Both $t$ and $\kappa$ are unchanged by scalar-unit cancellation. Moreover
$\kappa=0$ if and only if $t=0$. If $t>0$, then $1\le\kappa\le t$.
\end{proposition}

\begin{proof}
The homomorphism $R\to\mathbb F_2$ setting every $x_i$ to zero sends
every scalar-unit cancellation to a binary cancellation. At termination
its differential is zero, so its object count is $\dim H(V,A)=t$.
This is the existing residue-survivor invariant specialized to one matching.

The homomorphism $R\to A_\epsilon$ sending every $x_i$ to $\epsilon$
sends $Q$ to $A+\epsilon B$. In the basis $1,\epsilon$ its underlying
binary differential is the doubled matrix in
\eqref{eq:first-jet-two-ranks}; its homology dimension is $2\kappa$.
Scalar-unit cancellation remains a homotopy equivalence after substitution,
so $\kappa$ is invariant without any cobordism-valued Schur update.

Choose a binary special contraction of $(V,A)$ onto $H(V,A)$. Transferring
$\epsilon B$ retains only its first perturbation term because
$\epsilon^2=0$. The transferred differential is $\epsilon B_*$ and its
binary rank is $r(B_*)$. Its homology dimension is $2t-2r(B_*)$, giving
the asserted formula.

Regard $B_*$ as a representation of a finite linearly oriented quiver;
it need not square to zero. The number of intervals in its decomposition
is $t-\sum_h r(B_{*,h})=\kappa$. A nonzero bounded quiver has at least
one interval. If $t=0$, scalar-unit cancellation removes the whole
relative complex, so it is contractible.
\end{proof}

\begin{lemma}[Established knot-completion multiplier]
\label{lem:first-jet-knot-multiplier}
If the actual classical completion $L_m$ is a knot, then
\begin{equation}
 \dim H(\operatorname{cl}(C))
 =\kappa\,\dim\operatorname{Kh}(L_m;\mathbb F_2).
 \label{eq:first-jet-knot-multiplier}
\end{equation}
\end{lemma}

This is the closure formula already established in the repository's
report~28 and \path{synthesis/closure_resets.tex}. We reproduce its
chain-level justification to make the pre-elimination query self-contained~\cite{ClosurePrior}.
The new all-link extension at $\kappa=1$ is proved separately below.

\begin{proof}
Let $(D,d_D)$ be the ordinary Khovanov complex of $L_m$. Its marked-arc
dot maps $X_i$ commute with each other and with $d_D$, and satisfy
$X_i^2=0$. On $D^a\otimes V^b$ use total cohomological degree $a+b$.
The completed differential is
$\mathcal D=d_D\otimes1+Q(X_1,\ldots,X_p)$; both summands have total
degree one. All tensor signs disappear over $\mathbb F_2$.

Consider moving the mark for variable $x_i$ across one crossing, from
edge $u$ to the continuation edge $v$ of the same link strand.
The local reverse-saddle map $H$ has cohomological degree $-1$,
satisfies $H^2=0$, commutes with every fixed dot map, and obeys
\[
 d_DH+Hd_D=X_u+X_v=\delta,\qquad H\delta=\delta H=0.
\]
These are strict chain identities, not merely identities on homology:
Nahm records them in Lemma~2.4, with attribution to
Hedden--Ni, Baldwin--Levine--Sarkar, and Lipshitz--Sarkar
\cite{Nahm2026Basepoints}. The last identity also follows there by combining
$HX_u=X_vH$ with commutation of $H$ with $X_u$.

Put $P=\partial Q/\partial x_i$ and evaluate its other variables at the
current marked edges. It has degree one on $V$ and is independent of
$x_i$. Differentiating $Q^2=0$ gives $PQ+QP=0$. The map $HP$ on
$D\otimes V$ has total degree zero, and
$G=1+HP$ is self-inverse since $H^2=0$.
Expansion, with compositions acting right to left, gives
\[
 G\mathcal D G
 =\mathcal D+\delta P+Hd_DH\,P^2
 =\mathcal D+\delta P .
\]
Here the linear $Q$-terms cancel by $PQ+QP=0$, the quadratic
$Q$-terms contain $H^2$, and
$Hd_DH=H\delta=0$ kills the remaining quadratic term.
In particular no assumption $P^2=0$ is being made.
Since each coefficient is square-free in $x_i$, adding $\delta P$
replaces exactly the action $X_u$ by $X_v$ in that variable.
If quantum gradings are retained, $H$ has quantum degree $-2$ and
differentiation with respect to $x_i$ raises quantum degree by two,
so the same $G$ has quantum degree zero for homogeneous $Q$.

Because $L_m$ is a knot, every mark can be moved along its strand to
one chosen edge by finitely many such conjugations. Their composition
is a chain isomorphism, providing coherent transport of the entire
polynomial matrix. Once all actions equal $X$, every monomial of
degree at least two vanishes and the differential is $d_D+A+XB$.
An equality of dot maps on homology alone would not justify this step.

For completeness, let $\nu$ be the operator changing one circle label
$x$ to $1$ and summing over circles. Over $\mathbb F_2$,
\[
 \nu^2=0,\qquad \nu d_D=d_D\nu,\qquad \nu X+X\nu=1;
\]
see Shumakovitch, Lemma~3.1.C and the proof of Theorem~3.2.A
\cite{Shumakovitch2014}. Thus $E=\operatorname{im}\nu=\ker\nu$
is a subcomplex and $D=E\oplus XE\cong A_\epsilon\otimes E$,
with $X$ acting as multiplication by $\epsilon$.
The transported completed complex is therefore the tensor product
over $\mathbb F_2$ of $E$ and $(A_\epsilon\otimes V,A+\epsilon B)$.
The latter has homology dimension $2\kappa$ by
Proposition~\ref{prop:first-jet-invariants}; the former has homology
dimension $\frac12\dim H(D)$.
The field K\"unneth formula gives \eqref{eq:first-jet-knot-multiplier}.
\end{proof}

\subsection{Primitive-interval rigidity}

\begin{lemma}[A primitive radical element is a coordinate]
\label{lem:first-jet-coordinate}
If $\theta\in R$ has zero constant term and nonzero linear part, an
automorphism of $R$ carries $x_1$ to $\theta$. Consequently
\[
 \operatorname{ann}_R(\theta)=\theta R.
\]
If $\eta$ also has nonzero linear part and $\theta\eta=0$, then
$\eta=\theta u$ for some unit $u\in R$.
\end{lemma}

\begin{proof}
Every element of $\mathfrak m=(x_1,\ldots,x_p)$ squares to zero: cross
terms cancel in characteristic two and squares of nonconstant monomials
vanish. Complete the linear part of $\theta$ to a basis of
$\mathfrak m/\mathfrak m^2$, lift the other basis vectors, and send the
generators to these $p$ elements. Their squares vanish, so this defines
an endomorphism of $R$. The induced map on the associated graded ring
of the finite $\mathfrak m$-adic filtration is an invertible linear
change of generators. The endomorphism is therefore an automorphism.

Write $R=R'\oplus x_1R'$ for
$R'=\mathbb F_2[x_2,\ldots,x_p]/(x_i^2)$. Multiplication by $x_1$
kills exactly $x_1R'$, so $\operatorname{ann}(x_1)=x_1R$.
Transport this identity by the automorphism. Thus
$\theta\eta=0$ gives $\eta=\theta u$. If $u$ had zero constant term,
the product would have zero linear part. Its constant term must
therefore be one, so $u$ is a unit.
\end{proof}

\begin{theorem}[Rigidity at multiplier one]
\label{thm:first-jet-rigidity}
If $\kappa=1$, then after scalar-unit cancellation $C$ is isomorphic
over $R$ to one finite interval
\[
 R\xrightarrow{\theta}R\xrightarrow{\theta}\cdots
 \xrightarrow{\theta}R
\]
on $m$, with $\theta$ of odd total-linear parity. A one-term interval
is allowed and needs no coefficient. For every nonempty classical link
completion of $m$,
\begin{equation}
 \dim H(\operatorname{cl}(C))
 =\dim\operatorname{Kh}(L_m;\mathbb F_2).
 \label{eq:first-jet-all-link-equality}
\end{equation}
In particular no knot-completion premise is needed here. Quantum
homogeneity and genuine cube origin are not assumptions of this theorem.
\end{theorem}

\begin{proof}
After cancelling scalar units, the quiver of total-linear matrices has
exactly one interval by Proposition~\ref{prop:first-jet-invariants}.
Hence its spaces are one-dimensional in consecutive degrees and zero
elsewhere, and its intervening binary arrows are one.
The full differential has one coefficient $\theta_h$ between consecutive
objects, and each $\theta_h$ has odd total-linear parity.

The chain condition gives $\theta_h\theta_{h+1}=0$.
Lemma~\ref{lem:first-jet-coordinate} gives
$\theta_{h+1}=\theta_hu_h$ with $u_h$ a unit. Inductively each coefficient
is a unit multiple of the first one. Successive rescalings of the
rank-one free modules make every arrow equal to the same $\theta$.
This is a chain isomorphism over the full ring, not just its first jet.

Let $D$ be the unreduced characteristic-two Khovanov complex of the
classical link $L_m$, with commuting square-zero dot operators $X_i$.
The vertical operator $\nu$, which changes one circle label $x$ to $1$
and sums over circles, satisfies
\[
 \nu^2=0,\qquad \nu d_D=d_D\nu,\qquad \nu X_i+X_i\nu=1.
\]
These identities follow directly from multiplication and comultiplication
in $\mathbb F_2[x]/(x^2)$ and hold for links as well as knots. Put
$\Theta=\theta(X_1,\ldots,X_p)$ and
\[
 U=\nu\Theta+\Theta\nu
 =\left(\sum_i\frac{\partial\theta}{\partial x_i}\right)(X_1,\ldots,X_p).
\]
The polynomial defining $U$ has constant term one. Every zero-constant
polynomial in commuting square-zero operators squares to zero, so $U^2=1$.
The derivation $\partial=\sum_i\partial/\partial x_i$ satisfies
$\partial^2=0$, giving $\nu U=U\nu$.
All these operators commute with $d_D$. Thus $K=U\nu$ satisfies
\[
 K^2=0,\qquad K\Theta+\Theta K=1,\qquad Kd_D=d_DK.
\]
For $E=\operatorname{im}K=\ker K$, these relations give a direct sum
of subcomplexes
$D=E\oplus\Theta E\cong A_\epsilon\otimes E$, with $\Theta$ acting
as multiplication by $\epsilon$.
The completed $\ell$-term interval is the tensor product of $E$ with
the free dual-number interval $J_\ell$. As a binary complex $J_\ell$
has dimension $2\ell$, differential rank $\ell-1$, and homology
dimension two. The field K\"unneth formula gives
$\dim H(J_\ell\otimes E)=2\dim H(E)=\dim H(D)$.
\end{proof}

The two ranks certify existence of the interval decomposition and
rescaling; no such basis needs to be constructed by the decision query.
The replacement asserts neither an isotopy of knots nor equality of
open complexes or of original bigraded homology.

\subsection{The boundary of the link extension}

For $p=2$, the two-term complex $R^2\to R$ with row $(x_1,x_2)$ has
$t=3$ and $\kappa=2$. Completing to one circle identifies the variables
and gives homology dimension four, as the knot formula predicts.
For the two-component unlink the module is $R$ itself. The image is
$(x_1,x_2)$, which has dimension three. The total homology dimension is
$12-2\cdot3=6$, whereas
$\kappa\,\dim\operatorname{Kh}(U_2)=2\cdot4=8$.
Thus the all-link equality cannot be extended to arbitrary $\kappa$.
A three-term interval with coefficient $x_1+x_2$ similarly has
$\kappa=3$ but completed rank four on the two-component unlink.

Raw survivor count alone is not a general algebraic decision rule:
an interval with coefficient $x_1$ can have arbitrarily many objects
and still have $\kappa=1$. Such algebraic fixtures must not be
presented as realized knot prefixes.

\subsection{A flat-origin constraint on genuine minimal complexes}

\begin{proposition}[Vanishing of the minimal diagonal first jet]
\label{prop:first-jet-flat-origin}
Start from an ordinary undotted tangle cube over $\mathbb F_2$, apply
delooping and scalar-unit cancellation, and retain any whole relative
direct summand of the resulting minimal complex $M$. Every same-matching
entry of its differential has total-linear dot parity zero.
In particular a one-matching summand has $\kappa=t$.
\end{proposition}

\begin{proof}
In the dotted cobordism category, let $\partial$ remove one dot and sum
over all choices. It is a derivation for composition and respects the
relations in characteristic two. A one-dotted sphere differentiates
to an undotted sphere, hence to zero. The derivative of the two-dot
relation consists of two equal terms. The two dotted terms in
neck-cutting also differentiate to two equal undotted surfaces.
Their derivatives vanish. The undotted cube differential $d$ therefore
satisfies $\partial d=0$.

For a chain isomorphism changing $d$ to $D=FdF^{-1}$, differentiation
of the inverse identity gives
\[
 \partial D=\Gamma D+D\Gamma,\qquad \Gamma=(\partial F)F^{-1}.
\]
Delooping is such an isomorphism. Gaussian cancellation is a chain
isomorphism displaying a contractible disk as a direct summand;
retain these disks formally until the end. The block of this equation
on the remaining minimal summand still has the form
$\partial D_M=\Gamma_M D_M+D_M\Gamma_M$, because the displayed
differential has no off-block entries.

Recall the scalar-residue map $\rho$ from the maintained radical analysis
\cite{RadicalPrior}. At a frontier of $2k$ endpoints, a dotted basis
morphism $x_S:a\to b$ has nonnegative weight
$k-c(a,b)+2|S|$, where $c(a,b)$ counts the circles in the glued pair
of matchings. Weight zero consists precisely of scalar identities.
The Euler-characteristic composition rule makes weights additive in
every nonzero valid composition. Projection to weight zero is therefore
multiplicative; this is the map $\rho$.
Minimality gives $\rho(D_M)=0$, hence
$\rho(\partial D_M)=0$. On a same-matching entry this is exactly
the sum of its singleton-dot coefficients. The residue on a
different-matching entry is automatically zero. Restriction to a further
whole summand preserves the same block argument.
\end{proof}

This proposition needs genuine cube origin or a suitable connection
identity. The generic query does not infer it from $Q^2=0$ or from a
nonnegative survivor profile. It computes both ranks and remains valid
on the broader algebraic class in Theorem~\ref{thm:first-jet-rigidity}.
A future provenance-dependent optimization could omit the second rank;
this release does not install that shortcut.

\subsection{Decision control, costs, and the remaining structural question}

\begin{proposition}[Established rank-two unknot test over $\mathbb F_2$]
\label{prop:first-jet-rank-two}
For a classical knot $K$,
\[
 \dim_{\mathbb F_2}\operatorname{Kh}(K;\mathbb F_2)=2
 \quad\Longleftrightarrow\quad K\text{ is the unknot}.
\]
Every larger total unreduced rank is at least four.
\end{proposition}

\begin{proof}
Choose a basepoint and its dot map $X$ on the ordinary complex $D$.
The reduced complex can be taken as $\widetilde D=\operatorname{im}X$,
up to its conventional quantum shift. The identities for $\nu$ above give
$D=E\oplus XE$ for $E=\operatorname{im}\nu$, and
$X:E\to XE=\widetilde D$ is an isomorphism with inverse $\nu$.
Both maps commute with the homological differential. Consequently
\[
 \dim_{\mathbb F_2}\operatorname{Kh}(K;\mathbb F_2)
 =2\dim_{\mathbb F_2}\widetilde{\operatorname{Kh}}(K;\mathbb F_2).
\]
This is the characteristic-two reduced/unreduced splitting of
Shumakovitch, Corollary~3.2.C \cite{Shumakovitch2014}.

The integral reduced cube is a bounded complex of finite free abelian
groups. The universal coefficient theorem therefore gives
\[
 \dim_{\mathbb Q}\widetilde{\operatorname{Kh}}(K;\mathbb Q)
 \le
 \dim_{\mathbb F_2}\widetilde{\operatorname{Kh}}(K;\mathbb F_2).
\]
Kronheimer--Mrowka prove that the rational reduced rank is greater
than one for every nontrivial knot: this is the rank consequence of
their instanton spectral sequence and unknot-detection argument,
specifically Corollary~1.3 and Proposition~1.4 with the accompanying
discussion \cite{KM2011Kh}. Hence a nontrivial knot has unreduced
$\mathbb F_2$ rank at least four. Conversely the crossing-free unknot
has complex $\mathbb F_2[x]/(x^2)$ with zero differential and rank two.
All normalization shifts leave these total dimensions unchanged.
\end{proof}

This detection theorem is an established topological input.
Our contribution is the preservation of its total-rank observation
under the new queries and replacements, not a new proof of Khovanov
unknot detection.


For a classical link $L$ with $c\ge1$ components,
$\dim_{\mathbb F_2}\operatorname{Kh}(L;\mathbb F_2)\ge2^c$.
Indeed, write $J_L(q)$ for its graded Euler characteristic, normalized
by $J_U(q)=q+q^{-1}$. The integral cube is free, so this Euler
characteristic is unchanged by passage to $\mathbb F_2$.
The skein relation
$q^{-2}J_{L_+}-q^2J_{L_-}=(q^{-1}-q)J_{L_0}$
shows that crossing changes preserve $J_L(1)$.
Changing crossings to obtain the $c$-component unlink gives
$J_L(1)=2^c$, and total homology dimension bounds the absolute
Euler characteristic \cite{Khovanov2000}.
Homological or quantum normalization shifts do not alter this
absolute bound. This argument is applied to the ordinary classical
completion, whose total rank the replacement theorem preserves.

At each proper prefix the new driver fully allocates one crossing with
cancellation disabled and finds whole differential components.
A component with $\kappa=0$ is discarded; one with $\kappa=1$ can be
replaced by one copy of its matching, with an arbitrary homological shift
because only total rank is observed. Every nonempty classical completion
of the latter contributes at least two to rank, and a completion with
at least two link components contributes at least four.
For $\kappa\ge2$ the multiplier is used only after verifying that the
actual completion is a classical knot.
Contributions of independent components are added; a lower bound of
four rejects the original knot.

Every replacement preserves total rank for every later classical
completion. This statement is preserved by finite sequences of
replacements; the theorem does not require quantum homogeneity of the
intermediate rank-only model. The driver then uses ordinary scalar
cancellation and continues its crossing order. At final closure it
reports rank two as unknotted and larger rank as knotted, using
Khovanov unknot detection.

The local optional allowance covers component traversal, polynomial
extraction, binary work, completion geometry, and construction of
replacement arrays. Arrays and reverse indices are prepared before a
final check and one installation. Exhausting this allowance discards
all uninstalled proposals from that stage, disables further observation,
and resumes ordinary cancellation on the intact state. Global
resource exhaustion propagates the ordinary scan limit and never
becomes a knot verdict. The object ceiling includes temporary crossing
allocation before observation.

For a queried block with $N=\sum_h n_h$ objects and total packed
coefficient input length $L$, the query stores binary matrices using
$O(N^2)$ bits up to index overhead and performs no cobordism composition.
A conservative bit-complexity envelope is
\[
 \widetilde O(L+N^3).
\]
The tilde absorbs logarithmic factors and polynomial-size index
arithmetic. Singleton monomials occupy packed bit positions
$1,2,4,\ldots$; reading only these avoids enumeration of dense
irrelevant higher terms. The cost includes the given coefficient
representation and does not assume it has polynomial size in $p$.

Let $w$ be the maximum actual frontier size and $g$ the largest number
of crossings between successive one-object checkpoints of the executed
scan, including its initial and terminal segments. The crossing ending
a segment counts before observation or cancellation. A crossing creates
at most eight delooped children, so each segment has at most $8^g$
temporary objects. Coefficient spaces at width $w$ have size $2^{O(w)}$.
The explicit elimination, query, and geometry operations therefore give
\[
 T(n;g,w),\ S(n;g,w)\le\operatorname{poly}(n)2^{O(g+w)}
\]
with the maintained polynomial ordering overhead. A one-object
checkpoint can retain a large frontier, so $w$ cannot be dropped as
it can after a suitable diagram-splicing reset.
Quasi-polynomial time follows if constructible orders guarantee
$g+w=O(\log^2(n+2))$. No such guarantee for arbitrary knots or unknots
is established. Stages after observer exhaustion remain included in
the actual recorded gap.

\subsection{Finite evidence and interpretation}

The discovery run uses seed $29002$, accepting $180$ knot braid diagrams
from $591$ attempts. Greedy and shuffled orders give $360$ complete scans
and $2898$ stages. There are $117$ eligible pre-Schur one-matching
components, containing $633$ objects and $237$ scalar pairs.
All are nonsingleton and contain scalar units. There are $96$ cases
with $t=\kappa=1$ and $21$ with $t=\kappa\ge2$.
All $117$ pre/post-cancellation comparisons agree.
None of $8183$ inspected minimal same-matching entries has odd
total-linear parity, consistently with the flat-origin proposition.

These counts establish real eligibility before Schur updates, rather
than only synthetic interval fixtures. They do not establish a timing
gain; queries and traversal incur overhead. Paired timing data must be
read separately. In particular no $t>\kappa$ primitive-interval
compression occurs in this actual-diagram corpus.

The tests compare decisions with an independent complete crossing cube,
test continued scans after replacement and exhausted local budgets,
and evaluate polynomial complexes with scalar units and nonlinear terms
under random invertible changes of basis. They retain explicit
link-completion counterexamples and an invalid polynomial complex whose
first-jet identities still hold, documenting the precise audit scope.

\paragraph{Further research questions.}
Can a cheaply maintained cube-origin connection certificate eliminate
the second rank throughout the decision scanner? Can small scalar
components be queried without traversing the whole differential graph
while still certifying absence of external attachments?
Which geometric classes admit polylogarithmic one-object checkpoint
gaps and frontier widths? Can the flat-origin constraint be lifted beyond
scalar residue to restrict higher jets of mixed-matching minimal
complexes? Can hard reduced knot families be found where avoided
cobordism work outweighs observer overhead?

